% =========================================================================
% บทที่ 5: ทัศนศาสตร์เชิงคลื่น แสงเลเซอร์ และการแทรกสอดเชิงแสง
% =========================================================================

\chapter{ทัศนศาสตร์เชิงคลื่น แสงเลเซอร์ และการแทรกสอดเชิงแสง}
\label{ch:chapter05}

\noindent{\large\bfseries\color{physicsblue} ความอาพันธ์ของคลื่นแสง การแทรกสอดช่องแคบคู่ และการเลี้ยวเบนฟรอนโฮเฟอร์}

\vspace{0.4cm}

\begin{equationsbox}{สมการการแทรกสอดและการเลี้ยวเบนของคลื่นแสง}
\begin{align}
I(\theta) &= I_0 \cos^2\left(\frac{\pi d \sin\theta}{\lambda}\right) \left[\frac{\sin\left(\frac{\pi a \sin\theta}{\lambda}\right)}{\frac{\pi a \sin\theta}{\lambda}}\right]^2 \\
d\sin\theta &= m\lambda \text{ (Double-slit Maxima)}, \quad a\sin\theta = m'\lambda \text{ (Single-slit Minima)}
\end{align}
\end{equationsbox}

\section{ทฤษฎีและกรอบแนวคิดสำคัญ}

ทัศนศาสตร์เชิงคลื่น (Wave Optics หรือ Physical Optics) อธิบายปรากฏการณ์ของแสงที่ไม่สามารถอธิบายได้ด้วยทัศนศาสตร์เชิงเรขาคณิต เช่น การแทรกสอด (Interference) การเลี้ยวเบน (Diffraction) และโพลาไรเซชัน (Polarization) รากฐานสำคัญของการเกิดริ้วการแทรกสอดที่เสถียรคือ ความอาพันธ์ของคลื่นแสง (Optical Coherence) ทั้งความอาพันธ์เชิงเวลา (Temporal Coherence) ซึ่งสัมพันธ์กับความเป็นสีเดียว และความอาพันธ์เชิงพื้นที่ (Spatial Coherence) ซึ่งสัมพันธ์กับขนาดเชิงมุมของแหล่งกำเนิดแสง

การกำเนิดแสงเลเซอร์ (LASER: Light Amplification by Stimulated Emission of Radiation) อาศัยกระบวนการปล่อยรังสีแบบถูกกระตุ้น (Stimulated Emission) ร่วมกับการทำให้เกิดการผกผันของประชากร (Population Inversion) ภายในโพรงกำทอนเชิงแสง (Optical Resonator) ทำให้ลำแสงเลเซอร์มีเฟสตรงกันและมีความเข้มสว่างสูงยิ่ง

\section{ตัวอย่างการคำนวณตามกรอบวิเคราะห์ P-S-C-T}

\begin{psctexample}{5.1}{การวิเคราะห์ริ้วการแทรกสอดร่วมกับการเลี้ยวเบนในการทดลองช่องแคบคู่}
\textbf{1. ภาพและสถานการณ์ทางกายภาพ (Picture):}\\\\ แสงเลเซอร์ฮีเลียม-นีออน ($\lambda = 632.8\text{ nm}$) ตกกระทบช่องแคบคู่ที่มีระยะห่างระหว่างช่อง $d = 0.25\text{ mm}$ และความกว้างช่อง $a = 0.05\text{ mm}$ ฉากรับภาพอยู่ห่างออกไป $L = 2.0\text{ m}$

\vspace{4pt}
\textbf{2. กลยุทธ์และแบบจำลองทางฟิสิกส์ (Strategy):}\\\\ คำนวณจำนวนริ้วสว่างของการแทรกสอดที่ปรากฏอยู่ภายในแถบสว่างกลางของการเลี้ยวเบน (Central Diffraction Peak)

\vspace{4pt}
\textbf{3. ขั้นตอนการคำนวณเชิงตัวเลข (Calculation):}\\\\ แถบมืดแรกของการเลี้ยวเบนเกิดขึ้นที่มุม $\theta_1$ ซึ่งสอดคล้องกับ:
\begin{equation}
a \sin\theta_1 = \lambda \implies \sin\theta_1 = \frac{\lambda}{a}
\end{equation}
ริ้วสว่างของการแทรกสอดเกิดขึ้นที่มุม $\theta_m$ ซึ่ง:
\begin{equation}
d \sin\theta_m = m\lambda \implies \sin\theta_m = \frac{m\lambda}{d}
\end{equation}
ริ้วการแทรกสอดจะหายไป (Missing Orders) เมื่อตำแหน่งตรงกับแถบมืดของการเลี้ยวเบน:
\begin{equation}
m = \frac{d}{a} = \frac{0.25\text{ mm}}{0.05\text{ mm}} = 5
\end{equation}
ดังนั้น ริ้วที่ $m = \pm 5$ จะไม่ปรากฏ ริ้วสว่างของการแทรกสอดที่อยู่ในแถบสว่างกลางจึงมีค่า $m = 0, \pm 1, \pm 2, \pm 3, \pm 4$ รวมทั้งสิ้น $2(4) + 1 = 9$ ริ้ว

\vspace{4pt}
\textbf{4. การทดสอบความสมเหตุสมผล (Test):}\\\\ อัตราส่วน $d/a = 5$ กำหนดซองหุ้มการเลี้ยวเบน (Diffraction Envelope) ได้อย่างแม่นยำตรงกับผลการทดลองในห้องปฏิบัติการ
\end{psctexample}

\section{การเชื่อมโยงสู่การวิจัยฟิสิกส์ร่วมสมัย}

\begin{researchbox}{ข้อมูลเชิงลึกจากงานวิจัย}
การวัดการเลี้ยวเบนรังสีเอกซ์ (X-ray Diffraction: XRD) อาศัยกฎของแบร็กก์ ($2d\sin\theta = n\lambda$) ซึ่งเป็นหลักการเดียวกับทัศนศาสตร์เชิงคลื่น ในการวิเคราะห์โครงสร้างผลึกของสารสังเคราะห์และวัสดุชีวมวลคาร์บอน เพื่อตรวจสอบความเป็นผลึก (Crystallinity) และขนาดอนุภาคระดับนาโนเมตร
\end{researchbox}

\section{การฝึกแก้โจทย์ปัญหาเชิงระบบตามกรอบ C-C-A-F}

\begin{ccafsolution}{การวิเคราะห์เกณฑ์การจำแนกของเรย์ลีและกำลังแยกภาพของกล้องโทรทรรศน์}
\textbf{ขั้นที่ 1: การสร้างมโนทัศน์ (Conceptualize):}\\\\ กล้องโทรทรรศน์ออปติคัลขนาดเส้นผ่านศูนย์กลางปากกล้อง $D = 0.5\text{ m}$ ตรวจสังเกตดาวคู่ที่ความยาวคลื่น $\lambda = 550\text{ nm}$

\vspace{4pt}
\textbf{ขั้นที่ 2: การจำแนกประเภทโจทย์ (Categorize):}\\\\ ขีดจำกัดการเลี้ยวเบนและกำลังแยกภาพเชิงมุม (Diffraction Limit \& Angular Resolution)

\vspace{4pt}
\textbf{ขั้นที่ 3: การวิเคราะห์เชิงคณิตศาสตร์ (Analyze):}\\\\ ตามเกณฑ์ของเรย์ลี (Rayleigh's Criterion) สำหรับช่องเปิดวงกลม (Airy Disk):
\begin{equation}
\theta_{\text{min}} = 1.22 \frac{\lambda}{D} = 1.22 \frac{550 \times 10^{-9}\text{ m}}{0.5\text{ m}} = 1.342 \times 10^{-6}\text{ rad} \approx 0.277\text{ ฟิลิปดา}
\end{equation}
หากดาวคู่มีระยะห่างเชิงมุมน้อยกว่า $0.28''$ ภาพของวงแอรี่จะซ้อนทับกันจนไม่สามารถแยกออกเป็นดาวสองดวงได้

\vspace{4pt}
\textbf{ขั้นที่ 4: การสรุปและประเมินผล (Finalize):}\\\\ การเพิ่มขนาดหน้ากล้อง $D$ เป็นปัจจัยหลักที่ช่วยลดค่า $\theta_{\text{min}}$ และเพิ่มกำลังแยกภาพเชิงมุมของเครื่องมือวิจัยทางดาราศาสตร์
\end{ccafsolution}

\section{แบบฝึกหัดทบทวนและโจทย์ท้าทายประจำบท}

\begin{enumerate}[leftmargin=1.5cm, label={\textbf{\thechapter.\arabic*}}, itemsep=8pt]
    \item \textbf{โจทย์เชิงแนวคิด (Conceptual):} จงอธิบายความสำคัญทางกายภาพของกฎและสมการหลักในบทเรียนนี้ พร้อมยกตัวอย่างสถานการณ์จริงในธรรมชาติ
    \item \textbf{โจทย์การประมาณค่าเชิงฟิสิกส์ (Estimation):} จงใช้การวิเคราะห์เชิงมิติและอันดับของขนาด (Order of Magnitude) ในการประมาณค่าปริมาณสำคัญที่เกี่ยวข้อง
    \item \textbf{โจทย์วิจัยขั้นสูง (Research Challenge):} จงเขียนสคริปต์ Python จำลองพฤติกรรมของระบบตามสมการที่ได้ศึกษา พร้อมพลอตกราฟเปรียบเทียบผลลัพธ์เชิงทฤษฎี
\end{enumerate}

